\section{Wedderburn decomposition of the fixed-point algebra (continued)}

\begin{theorem}[Isotypic-component decomposition]%
    \label{thm:starSubalgebra_orthogonal_isotypic_decomposition}
    \lean{StarSubalgebra.exists_orthogonal_isotypic_decomposition}
    \leanok
    \uses{thm:starSubalgebra_orthogonal_irreducible_decomposition,
        lem:starSubalgebra_isOrtho_sSup_not_sameIsotype, lem:starSubalgebra_sameIsotype_of_le_sSup,
        lem:starSubalgebra_sameIsotype_refl, lem:starSubalgebra_sameIsotype_symm, lem:starSubalgebra_sameIsotype_trans}
    Let $S$ be a $*$-subalgebra of $\MN{D}$. Then $\C^D$ is the supremum of a finite family $\mathcal{C}$ of pairwise orthogonal subspaces, the
    isotypic components, and each component $C\in\mathcal{C}$ is the supremum $C=\bigvee_{W\in\mathcal{D}_C}W$ of a finite, nonempty class
    $\mathcal{D}_C$ of pairwise orthogonal subspaces irreducible under $S$ that are pairwise of the same type. Moreover, for distinct components
    $C\neq C'$, no subspace irreducible under $S$ contained in $C$ is of the same type as any subspace irreducible under $S$ contained in $C'$.
\end{theorem}

\begin{proof}
    \leanok
    \uses{thm:starSubalgebra_orthogonal_irreducible_decomposition,
        lem:starSubalgebra_isOrtho_sSup_not_sameIsotype, lem:starSubalgebra_sameIsotype_of_le_sSup,
        lem:starSubalgebra_sameIsotype_refl, lem:starSubalgebra_sameIsotype_symm, lem:starSubalgebra_sameIsotype_trans}
    Theorem~\ref{thm:starSubalgebra_orthogonal_irreducible_decomposition} supplies a finite, pairwise orthogonal family $\mathcal{D}$ of
    irreducible pieces with $\bigvee_{W\in\mathcal{D}}W=\C^D$. For each piece $W$, let $\mathcal{D}_W$ be the class of pieces of $\mathcal{D}$
    of the same type as $W$, and let $C_W$ be its supremum. Each piece lies in its own class by Lemma~\ref{lem:starSubalgebra_sameIsotype_refl}, so
    each class is nonempty, the classes cover $\mathcal{D}$, and the components span $\C^D$. Each class is contained in $\mathcal{D}$, so
    its pieces are pairwise orthogonal. Two pieces in one class are of the same type through the representative, by
    Lemmas~\ref{lem:starSubalgebra_sameIsotype_symm} and~\ref{lem:starSubalgebra_sameIsotype_trans}. If two representatives
    are of the same type, transitivity makes their classes, and hence their components, equal.

    Thus distinct components arise from representatives $W,W'$ that are not of the same type. Write $\sim$ for the same-type relation. For
    $a\in\mathcal{D}_W$ and $b\in\mathcal{D}_{W'}$, one has $W\sim a$ and $W'\sim b$. If $a\sim b$, symmetry and transitivity would give
    $W\sim a\sim b\sim W'$, a contradiction. Hence no piece of $\mathcal{D}_W$ is of the same type as a piece of
    $\mathcal{D}_{W'}$, and Lemma~\ref{lem:starSubalgebra_isOrtho_sSup_not_sameIsotype} makes the
    components orthogonal.

    Finally, let $V\leq C_W$ and $V'\leq C_{W'}$ be irreducible under $S$, with $C_W\neq C_{W'}$. By
    Lemma~\ref{lem:starSubalgebra_sameIsotype_of_le_sSup}, there are pieces $a\in\mathcal{D}_W$ and $b\in\mathcal{D}_{W'}$ with $V\sim a$ and
    $V'\sim b$. If $V\sim V'$, symmetry and transitivity would give $W\sim a\sim V\sim V'\sim b\sim W'$, again contradicting that $W$ and
    $W'$ are not of the same type.
\end{proof}

The pieces of one class are unitarily identified copies of a single irreducible piece. An orthonormal basis of one piece therefore spreads,
through the inner-product-preserving intertwiners, to an adapted orthonormal basis of the whole component. In this basis every member of $S$ acts by one
matrix on the irreducible index, identically across the multiplicity index. This is the tensor factorization of the isotypic components behind the block
form $\MN{d_k}\otimes\Id_{m_k}$ of
\cite[Theorem~6.14]{Wolf2012Quantum}.

\begin{theorem}[Adapted orthonormal basis of an isotypic component]%
    \label{thm:starSubalgebra_adapted_orthonormal_basis}
    \lean{StarSubalgebra.exists_adapted_orthonormal_family_of_sameIsotype}
    \leanok
    \uses{def:starSubalgebra_irreducibleSubspace, def:starSubalgebra_sameIsotype}
    Let $S$ be a $*$-subalgebra of $\MN{D}$ and let $\mathcal{D}$ be a finite, nonempty family of pairwise orthogonal subspaces of $\C^D$,
    irreducible under $S$ and pairwise of the same type. Write $m$ for the number of pieces in $\mathcal{D}$; all pieces share one dimension $d$.
    Then the component $C=\bigvee_{W\in\mathcal{D}}W$ has an orthonormal basis $(f_{i,j})_{0\leq i<m,\,0\leq j<d}$ such that for every $A\in S$
    there is a matrix $B\in\MN{d}$ satisfying
    \begin{align}
        A f_{i,j}
        &=\sum_{j'=0}^{d-1}B_{j'j}f_{i,j'}
        \label{eq:spectral_adapted_action}
    \end{align}
    for all $i<m$ and $j<d$. In the adapted basis, $A$ acts on $C$ by the matrix $B$ on the irreducible index, identically across the multiplicity
    index; after reordering the indices this is the block $\MN{d}\otimes\Id_m$ of~\cite[Theorem~6.14]{Wolf2012Quantum}. The
    vectors $(f_{i,j})_{0\leq j<d}$ of the $i$-th copy span one of the pieces of $\mathcal{D}$.
\end{theorem}

\begin{proof}
    \leanok
    \uses{lem:starSubalgebra_sameIsotype_isometry, lem:starSubalgebra_sameIsotype_refl}
    Fix a piece $W\in\mathcal{D}$, of dimension $d$, with an orthonormal basis $e_0,\ldots,e_{d-1}$, and enumerate
    $\mathcal{D}=\{W'_0,\ldots,W'_{m-1}\}$. Every piece of $\mathcal{D}$ is of the same type as $W$, by hypothesis for the other pieces and by
    Lemma~\ref{lem:starSubalgebra_sameIsotype_refl} for $W$ itself. Lemma~\ref{lem:starSubalgebra_sameIsotype_isometry} therefore provides
    intertwiners $u_i$ mapping $W$ onto $W'_i$ and preserving the inner product on $W$. Set $f_{i,j}=u_i(e_j)$.

    Within one copy, the vectors $f_{i,0},\ldots,f_{i,d-1}$ are orthonormal because $u_i$ preserves inner products; across copies they are
    orthogonal because distinct pieces of $\mathcal{D}$ are. Since $u_i$ maps $W$ onto $W'_i$, the vectors of the $i$-th copy span $W'_i$. Thus
    each $W'_i$ has dimension $d$, and all the vectors together span $C$.

    For $A\in S$, invariance of $W$ gives $Ae_j=\sum_{j'}B_{j'j}e_{j'}$, where
    $B_{j'j}=\langle e_{j'},Ae_j\rangle$. The intertwining identity gives
    \begin{align}
        A f_{i,j}
        &=A\,u_i(e_j)
          =u_i(Ae_j)
          =\sum_{j'}B_{j'j}u_i(e_{j'})
          =\sum_{j'}B_{j'j}f_{i,j'}.
        \notag
    \end{align}
    This is (\ref{eq:spectral_adapted_action}).
\end{proof}

\begin{theorem}[Spatial tensor factorization of the isotypic components]%
    \label{thm:starSubalgebra_isotypic_tensor_factorization}
    \lean{StarSubalgebra.exists_isotypic_tensor_factorization}
    \leanok
    \uses{thm:starSubalgebra_orthogonal_isotypic_decomposition,
        thm:starSubalgebra_adapted_orthonormal_basis} Let $S$ be a $*$-subalgebra of $\MN{D}$. Then $\C^D$ is the supremum
    of a finite family of pairwise orthogonal isotypic components, and each component carries an orthonormal basis $(f_{i,j})$, indexed by a
    multiplicity index $i<m$ and an irreducible index $j<d$, such that every $A\in S$ acts by
    \begin{align}
        A f_{i,j}
        &=\sum_{j'=0}^{d-1}B_{j'j}f_{i,j'}
        \notag
    \end{align}
    for a matrix $B\in\MN{d}$ depending only on $A$ and the component. In the adapted bases, the action of $S$ on each component is by
    matrix-times-identity blocks, the form $\MN{d_k}\otimes\Id_{m_k}$ of~\cite[Theorem~6.14]{Wolf2012Quantum}.
\end{theorem}

\begin{proof}
    \leanok
    \uses{thm:starSubalgebra_orthogonal_isotypic_decomposition,
        thm:starSubalgebra_adapted_orthonormal_basis} Theorem~\ref{thm:starSubalgebra_orthogonal_isotypic_decomposition}
    writes $\C^D$ as the supremum of finitely many pairwise orthogonal components, each the supremum of a finite, nonempty class of pairwise
    orthogonal irreducible pieces of one type. Applying Theorem~\ref{thm:starSubalgebra_adapted_orthonormal_basis} to each class
    equips each component with its adapted orthonormal basis.
\end{proof}

The adapted orthonormal bases of the pairwise orthogonal isotypic components concatenate into a single orthonormal basis of $\C^D$, indexed by a
component index $k$, a multiplicity index $i$, and an irreducible index $j$. In this basis every member of $S$ acts block-diagonally, one
matrix-times-identity block per component. A $*$-subalgebra contains the identity matrix, so every component carries irreducible pieces of $S$. This
is the unital case of the block representation invoked by~\cite[Theorem~6.14]{Wolf2012Quantum}, without the zero block.

\begin{theorem}[Adapted orthonormal basis of the whole space]%
    \label{thm:starSubalgebra_global_adapted_orthonormal_basis}
    \lean{StarSubalgebra.exists_adapted_orthonormalBasis}
    \leanok
    \uses{def:starSubalgebra_irreducibleSubspace, def:starSubalgebra_sameIsotype}
    Let $S$ be a $*$-subalgebra of $\MN{D}$. There are $K\in\N$, positive dimensions $d_0,\ldots,d_{K-1}$ and multiplicities
    $m_0,\ldots,m_{K-1}$, and an orthonormal basis $(f_{k,i,j})_{k<K,\,i<m_k,\,j<d_k}$ of $\C^D$ such that for every
    $A\in S$ there are matrices $B_k\in\MN{d_k}$ satisfying
    \begin{align}
        A f_{k,i,j}
        &=\sum_{j'=0}^{d_k-1}(B_k)_{j'j}f_{k,i,j'}
        \label{eq:spectral_global_action}
    \end{align}
    for all $k<K$, $i<m_k$, and $j<d_k$. In this basis, $A$ acts on the $k$-th component as $\Id_{m_k}\otimes B_k$. Each row
    $(f_{k,i,j})_{0\leq j<d_k}$ spans a subspace irreducible under $S$, and rows drawn from distinct components span subspaces that are never of the
    same type. Only the containment direction is asserted: every member of $S$ takes this form. Equality with the block algebra, including the
    reverse inclusion, is
    Theorem~\ref{thm:starSubalgebra_unitary_block_diagonal_iff}.
\end{theorem}

\begin{proof}
    \leanok
    \uses{thm:starSubalgebra_orthogonal_isotypic_decomposition,
        thm:starSubalgebra_adapted_orthonormal_basis}
    Theorem~\ref{thm:starSubalgebra_orthogonal_isotypic_decomposition}
    writes $\C^D$ as the supremum of finitely many pairwise orthogonal
    isotypic components, each the supremum of a finite, nonempty same-type
    class of pairwise orthogonal irreducible pieces.
    Theorem~\ref{thm:starSubalgebra_adapted_orthonormal_basis} equips the
    $k$-th component with an adapted orthonormal family
    $(f_{k,i,j})_{i<m_k,\,j<d_k}$ spanning it. Here $m_k$ is the number of
    pieces in the class, which is positive because the class is nonempty,
    and $d_k$ is their common dimension, which is positive because
    irreducible pieces are nonzero.

    Concatenating these families over $k$ yields an orthonormal family:
    vectors in one component are orthonormal by construction, and vectors
    in distinct components are orthogonal. The concatenated family spans
    $\C^D$ because each component is spanned by its family and the
    components have supremum $\C^D$; hence it is an orthonormal basis.
    The action property (\ref{eq:spectral_global_action}) holds
    componentwise. Each row spans one of the irreducible pieces in its
    component's class, and
    Theorem~\ref{thm:starSubalgebra_orthogonal_isotypic_decomposition}
    shows that irreducible subspaces in distinct components are never of the
    same type.
\end{proof}

The change of basis to an orthonormal basis indexed by such triples is
implemented by a unitary matrix, and it carries every matrix acting
blockwise on the basis to a block-diagonal matrix.

\begin{lemma}[Unitary change of basis to block-diagonal form]%
    \label{lem:starSubalgebra_unitary_change_of_basis_block}
    \lean{OrthonormalBasis.exists_unitary_conj_blockDiagonal}
    \leanok
    Let $(f_{k,i,j})_{k<K,\,i<m_k,\,j<d_k}$ be an orthonormal basis of
    $\C^D$. There are an identification of the index set of triples
    $(k,i,j)$ with $\{0,\ldots,D-1\}$, realizing
    $\sum_kd_km_k=D$, and a unitary $U\in\MN{D}$ whose columns are the
    basis vectors, such that every matrix $A\in\MN{D}$ acting on the basis
    by
    \begin{align}
        A f_{k,i,j}
        &=\sum_{j'=0}^{d_k-1}(B_k)_{j'j}f_{k,i,j'}
        \notag
    \end{align}
    for matrices $B_k\in\MN{d_k}$ satisfies
    \begin{align}
        U^\dagger A U
        &=\bigoplus_{k=0}^{K-1}\Id_{m_k}\otimes B_k.
        \label{eq:spectral_basis_block}
    \end{align}
\end{lemma}

\begin{proof}
    \leanok
    Let $U$ be the change-of-basis matrix from the standard orthonormal
    basis of $\C^D$ to $(f_{k,i,j})$, so the columns of $U$ are the basis
    vectors. A change of basis between orthonormal bases is unitary. Both
    bases have $D$ elements, which identifies the index set of triples with
    $\{0,\ldots,D-1\}$ and gives $\sum_kd_km_k=D$. The matrix
    $U^\dagger AU=U^{-1}AU$ represents $x\mapsto Ax$ in the basis
    $(f_{k,i,j})$, with entries
    \begin{align}
        \langle f_{k',i',j'},Af_{k,i,j}\rangle
        &=\sum_{j''}(B_k)_{j''j}
          \langle f_{k',i',j'},f_{k,i,j''}\rangle
          =\delta_{k'k}\delta_{i'i}(B_k)_{j'j}.
        \notag
    \end{align}
    This is the block-diagonal matrix in
    (\ref{eq:spectral_basis_block}).
\end{proof}

\begin{theorem}[Global unitary block-diagonal form]%
    \label{thm:starSubalgebra_unitary_block_diagonal}
    \lean{StarSubalgebra.exists_unitary_conj_blockDiagonal}
    \leanok
    Let $S$ be a $*$-subalgebra of $\MN{D}$. There are $K\in\N$, positive
    dimensions $d_0,\ldots,d_{K-1}$ and multiplicities
    $m_0,\ldots,m_{K-1}$ with $\sum_kd_km_k=D$, realized by an explicit
    identification of the index sets, and a unitary $U\in\MN{D}$ such that
    every $A\in S$ satisfies
    \begin{align}
        U^\dagger A U
        &=\bigoplus_{k=0}^{K-1}\Id_{m_k}\otimes B_k
        \label{eq:spectral_subalg_containment}
    \end{align}
    for matrices $B_k\in\MN{d_k}$ depending on $A$. Up to reordering the
    two tensor factors of each block, this is the containment direction of
    the unital case of the block representation of
    \cite[Theorem~6.14]{Wolf2012Quantum}: $S$ is carried by $U$ into the
    block algebra $\bigoplus_k\Id_{m_k}\otimes\MN{d_k}$. Equality, including
    the reverse inclusion, is
    Theorem~\ref{thm:starSubalgebra_unitary_block_diagonal_iff}.
\end{theorem}

\begin{proof}
    \leanok
    \uses{thm:starSubalgebra_global_adapted_orthonormal_basis,
        lem:starSubalgebra_unitary_change_of_basis_block}
    Let $(f_{k,i,j})$ be the adapted orthonormal basis of
    Theorem~\ref{thm:starSubalgebra_global_adapted_orthonormal_basis}, and
    let $U$ be the unitary supplied by
    Lemma~\ref{lem:starSubalgebra_unitary_change_of_basis_block}. Every
    $A\in S$ acts on the basis by matrices $B_k$ on the irreducible index,
    identically across the multiplicity index, so
    (\ref{eq:spectral_basis_block}) gives
    (\ref{eq:spectral_subalg_containment}).
\end{proof}

The reverse inclusion of the block representation rests on the
double-commutant characterization: in finite dimensions a $*$-subalgebra
contains every matrix that commutes with its commutant. The proof is the
Jacobson density theorem for $\C^D$ as a semisimple module over the
subalgebra.

\begin{lemma}[Membership through the double commutant]%
    \label{lem:starSubalgebra_double_commutant}
    \lean{StarSubalgebra.mem_of_forall_comm}
    \leanok
    Let $S$ be a $*$-subalgebra of $\MN{D}$ and let $T\in\MN{D}$. If $T$
    commutes with every linear operator on $\C^D$ that commutes with all
    members of $S$, then $T\in S$.
\end{lemma}

\begin{proof}
    \leanok
    \uses{lem:starSubalgebra_isSemisimpleModule}
    The space $\C^D$ is a semisimple module over $S$ by
    Lemma~\ref{lem:starSubalgebra_isSemisimpleModule}, and the linear
    operators commuting with all members of $S$ are exactly the
    endomorphisms of this module. The hypothesis therefore makes $T$ linear
    over the endomorphism ring of the module. By the Jacobson density
    theorem, for every finite set of vectors there is a member of $S$ acting
    on them as $T$ does. Applied to a basis of $\C^D$, this produces a
    member of $S$ whose action agrees with $T$ on a basis, hence equals $T$.
\end{proof}

\begin{theorem}[Complementary spatial actions of a $*$-subalgebra and its commutant]%
    \label{thm:starSubalgebra_complementary_spatial_actions}
    \lean{StarSubalgebra.exists_complementary_action_orthonormalBasis}
    \leanok
    Let $S$ be a $*$-subalgebra of $\MN{D}$. There are positive integers
    $d_k,m_k$ and an orthonormal basis
    $(f_{k,i,j})_{k<K,\,i<m_k,\,j<d_k}$ of $\C^D$ such that every $A\in S$
    acts by
    \begin{align}
        A f_{k,i,j}
        &=\sum_{j'=0}^{d_k-1}(B_k)_{j'j}f_{k,i,j'},
        \label{eq:spec_subalgebra_action}
    \end{align}
    while every $T\in\MN{D}$ commuting with all members of $S$ acts by
    \begin{align}
        T f_{k,i,j}
        &=\sum_{i'=0}^{m_k-1}(C_k)_{i'i}f_{k,i',j}.
        \label{eq:spec_commutant_action}
    \end{align}
    Thus the same orthonormal identification realizes $S$ on the second
    tensor factor and its commutant on the complementary first tensor
    factor. This is the finite-dimensional $C^*$-algebra step used in the
    proof of~\cite[Lemma~2.1]{Beigi2012Classification}.
\end{theorem}

\begin{proof}
    \leanok
    \uses{thm:starSubalgebra_global_adapted_orthonormal_basis}
    Choose the adapted orthonormal basis $(f_{k,i,j})$ for $S$.
    Formula (\ref{eq:spec_subalgebra_action}) is its defining action
    property. If $T$ commutes with $S$, the row transport
    \begin{align}
        \tau_{k,i',i}(v)
        &:=\sum_j\langle f_{k,i',j},v\rangle f_{k,i,j}
        \notag
    \end{align}
    also commutes with $S$, and so does $\tau_{k,i',i}T$. Schur's lemma
    makes this composite scalar on each irreducible row and makes its matrix
    coefficients vanish between distinct isotypic components. Expanding
    $Tf_{k,i,j}$ in the orthonormal basis therefore gives
    (\ref{eq:spec_commutant_action}), with coefficients
    $(C_k)_{i'i}$ independent of $j$.
\end{proof}

Combining the containment direction with the double-commutant criterion
yields the full block representation of a finite-dimensional
$*$-subalgebra, the unital case of Equation~(1.39) of
\cite{Wolf2012Quantum} invoked by
\cite[Theorem~6.14]{Wolf2012Quantum}: up to a unitary change of basis, $S$
equals the block algebra $\bigoplus_k\Id_{m_k}\otimes\MN{d_k}$.

\begin{theorem}[Full block form of a $*$-subalgebra]%
    \label{thm:starSubalgebra_unitary_block_diagonal_iff}
    \lean{StarSubalgebra.exists_unitary_conj_blockDiagonal_iff}
    \leanok
    Let $S$ be a $*$-subalgebra of $\MN{D}$. There are $K\in\N$, positive
    dimensions $d_0,\ldots,d_{K-1}$ and multiplicities
    $m_0,\ldots,m_{K-1}$ with $\sum_kd_km_k=D$, realized by an explicit
    identification of the index sets, and a unitary $U\in\MN{D}$ such that
    a matrix $A\in\MN{D}$ belongs to $S$ exactly when
    \begin{align}
        U^\dagger A U
        &=\bigoplus_{k=0}^{K-1}\Id_{m_k}\otimes B_k
        \label{eq:spectral_full_block}
    \end{align}
    for some matrices $B_k\in\MN{d_k}$; that is, up to reordering the two
    tensor factors of each block,
    \begin{align}
        S
        &=U\left(\bigoplus_{k=0}^{K-1}
          \Id_{m_k}\otimes\MN{d_k}\right)U^\dagger.
        \notag
    \end{align}
    This is the block representation of a finite-dimensional $*$-algebra in
    Equation~(1.39) of~\cite{Wolf2012Quantum}, invoked by
    \cite[Theorem~6.14]{Wolf2012Quantum}, in the unital case: a
    $*$-subalgebra contains the identity matrix, so there is no zero block.
\end{theorem}

\begin{proof}
    \leanok
    \uses{thm:starSubalgebra_global_adapted_orthonormal_basis,
        lem:starSubalgebra_unitary_change_of_basis_block,
        lem:starSubalgebra_double_commutant,
        lem:starSubalgebra_scalar_commutant_irreducible, def:starSubalgebra_intertwiner,
        def:starSubalgebra_sameIsotype}
    Let $(f_{k,i,j})$ be the adapted orthonormal basis of
    Theorem~\ref{thm:starSubalgebra_global_adapted_orthonormal_basis}, write
    $W_{k,i}$ for the span of the row $(f_{k,i,j})_{j<d_k}$, and let $U$ be
    the unitary supplied by
    Lemma~\ref{lem:starSubalgebra_unitary_change_of_basis_block}. Every
    member of $S$ satisfies (\ref{eq:spectral_full_block}), which is the
    containment direction.

    For the reverse inclusion, fix blocks $(B_k)_k$ and let $T$ act on the
    adapted basis by
    $Tf_{k,i,j}=\sum_{j'}(B_k)_{j'j}f_{k,i,j'}$. By
    Lemma~\ref{lem:starSubalgebra_double_commutant}, it suffices to show that
    $T$ commutes with every operator $g$ that commutes with all members of
    $S$. For a component index $k$ and multiplicity indices $i,i'$, consider
    \begin{align}
        \tau(v)
        &:=\sum_j\langle f_{k,i',j},v\rangle f_{k,i,j}.
        \notag
    \end{align}
    Because every member of $S$ acts on the rows $(k,i')$ and $(k,i)$ by
    the same matrix, $\tau$ commutes with all members of $S$, and so does
    $\tau\circ g$. The composite maps $\C^D$ into $W_{k,i}$, and
    $W_{k,i}$ is irreducible. Hence
    Lemma~\ref{lem:starSubalgebra_scalar_commutant_irreducible} gives a
    scalar $\gamma^{(k)}_{i'i}\in\C$ such that
    $\tau(gx)=\gamma^{(k)}_{i'i}x$ for every $x\in W_{k,i}$. Evaluating at
    $x=f_{k,i,j}$ gives
    \begin{align}
        \langle f_{k,i',j'},g f_{k,i,j}\rangle
        &=\gamma^{(k)}_{i'i}\delta_{j'j}.
        \label{eq:spectral_commutant_coeff}
    \end{align}

    Matrix coefficients of $g$ between distinct components vanish. Indeed,
    for $k'\neq k$, define $\tau$ using the row $(k',i')$. Then
    $\tau\circ g$ is an intertwiner from $W_{k,i}$ into $W_{k',i'}$. A
    nonzero coefficient would make this intertwiner nonzero on $W_{k,i}$,
    so the two rows would be of the same type, contrary to
    Theorem~\ref{thm:starSubalgebra_global_adapted_orthonormal_basis}.
    Expanding $gf_{k,i,j}$ and using
    (\ref{eq:spectral_commutant_coeff}) therefore gives
    \begin{align}
        g f_{k,i,j}
        &=\sum_{i'}\gamma^{(k)}_{i'i}f_{k,i',j}.
        \notag
    \end{align}
    Thus $g$ has blocks $C_k\otimes\Id_{d_k}$ in the adapted basis,
    complementary to the blocks of $T$. The two actions commute on every
    basis vector, so $Tg=gT$. By
    Lemma~\ref{lem:starSubalgebra_double_commutant}, $T\in S$.

    Finally, if $A$ satisfies (\ref{eq:spectral_full_block}), then
    $U^\dagger AU=U^\dagger TU$ for the member $T\in S$ constructed from
    the blocks $(B_k)_k$. Hence $A=T\in S$.
\end{proof}

\begin{theorem}[Fixed-point algebra is semisimple]%
    \label{thm:fixedPointAlgebra_isSemisimpleRing}
    \lean{Kraus.fixedPointAlgebra_isSemisimpleRing}
    \leanok
    \uses{thm:adjointFixedPointsStarSubalgebra}
    The adjoint-fixed-point $*$-subalgebra of a trace-preserving Kraus map
    on $\MN{D}$ is a semisimple ring.
\end{theorem}

\begin{proof}
    \leanok
    \uses{thm:adjointFixedPointsStarSubalgebra}
    The subalgebra is finite-dimensional as a subalgebra of $\MN{D}$, hence
    Artinian. Following~\cite[Theorem~6.14]{Wolf2012Quantum}, it suffices to
    show that the Jacobson radical $J$ vanishes. If $x\in J$, then
    $x^*x\in J$ by left-ideal closure. The radical of an Artinian ring is
    nilpotent, so $x^*x$ is nilpotent. A self-adjoint nilpotent matrix
    satisfies
    $\lVert x^*x\rVert^{2^k}=\lVert(x^*x)^{2^k}\rVert=0$, hence
    $x^*x=0$, and the $C^*$-identity gives $x=0$.
\end{proof}

\begin{theorem}[Abstract Wedderburn--Artin decomposition]%
    \label{thm:fixedPointAlgebra_wedderburnArtin}
    \lean{Kraus.fixedPointAlgebra_wedderburnArtin}
    \leanok
    \uses{thm:fixedPointAlgebra_isSemisimpleRing}
    There exist $n\in\N$ and dimensions $d_1,\ldots,d_n\geq1$ such that
    the adjoint-fixed-point $*$-subalgebra is $\C$-algebra isomorphic to
    $\prod_{k=1}^{n}\MN{d_k}$.
\end{theorem}

\begin{proof}
    \leanok
    \uses{thm:fixedPointAlgebra_isSemisimpleRing}
    Combine semisimplicity
    (Theorem~\ref{thm:fixedPointAlgebra_isSemisimpleRing}) with the
    Wedderburn--Artin structure theorem for finite-dimensional semisimple
    algebras over an algebraically closed field.
\end{proof}

\begin{theorem}[Wedderburn--Artin decomposition of the unital fixed-point algebra]%
    \label{thm:fixedPointsStarSubalgebra_exists_algEquiv_pi_matrix}
    \lean{Kraus.fixedPointsStarSubalgebra_exists_algEquiv_pi_matrix}
    \leanok
    \uses{thm:starSubalgebra_wedderburnArtin, thm:fixedPointsStarSubalgebra}
    Let $E$ be a unital Kraus map on $\MN{D}$ whose adjoint map $E^*$ has a
    positive definite fixed point $\rho>0$. There exist $n\in\N$ and
    dimensions $d_1,\ldots,d_n\geq1$ such that the fixed-point
    $*$-subalgebra is $\C$-algebra isomorphic to
    $\Fix(E)\cong\prod_{k=1}^{n}\MN{d_k}$.
\end{theorem}

\begin{proof}
    \leanok
    \uses{thm:starSubalgebra_wedderburnArtin, thm:fixedPointsStarSubalgebra}
    By Theorem~\ref{thm:fixedPointsStarSubalgebra}, the fixed-point set
    $\Fix(E)$ is a $*$-subalgebra of $\MN{D}$. Applying
    Theorem~\ref{thm:starSubalgebra_wedderburnArtin} gives the asserted
    product decomposition with each block size $d_k\geq1$.
\end{proof}

\begin{theorem}[Fixed-point algebra decomposition with multiplicities]%
    \label{thm:adjointFixedPoints_wedderburnDecomp}
    \lean{Kraus.adjointFixedPoints_wedderburnDecomp}
    \leanok
    \uses{thm:fixedPointAlgebra_wedderburnArtin}
    There exist integers $n$, block sizes $d_1,\ldots,d_n\geq1$,
    multiplicities $m_1,\ldots,m_n\geq1$, and a $\C$-algebra isomorphism
    from the adjoint-fixed-point $*$-subalgebra to
    $\prod_{k=1}^{n}\MN{d_k}$ such that
    $\sum_{k=1}^{n}d_km_k\leq D$.
\end{theorem}

\begin{proof}
    \leanok
    \uses{thm:fixedPointAlgebra_wedderburnArtin}
    Apply the abstract Wedderburn decomposition
    (Theorem~\ref{thm:fixedPointAlgebra_wedderburnArtin}). Since the
    adjoint-fixed-point algebra is realized as a subalgebra of $\MN{D}$,
    this realization also yields multiplicities $m_k$ and the bound
    $\sum_kd_km_k\leq D$.
\end{proof}

\begin{remark}[Block form of the fixed-point algebra]
    Wolf's block form identifies the adjoint-fixed-point algebra, up to a
    unitary change of basis, with
    $0\oplus\bigoplus_k\MN{d_k}\otimes\Id_{m_k}$, together with the
    density-block form $\MN{d_k}\otimes\rho_k$.
    Theorem~\ref{thm:adjointFixedPoints_wedderburnDecomp} gives the product
    decomposition, multiplicities, and dimension bound from this block
    form; the unitary conjugation identity itself, in the unital case, is
    Theorem~\ref{thm:adjointFixedPoints_block_form}.
\end{remark}

\begin{theorem}[Unitary block form of the Heisenberg-picture fixed points]%
    \label{thm:adjointFixedPoints_block_form}
    \lean{Kraus.adjointFixedPoints_blockDiagonal_iff}
    \leanok
    \uses{def:adjoint_fixed_points, def:tp_kraus}
    Let $E$ be a trace-preserving Kraus map on $\MN{D}$ with a positive
    definite fixed point $\rho>0$, $E(\rho)=\rho$. There are $n\in\N$,
    positive dimensions $d_0,\ldots,d_{n-1}$ and multiplicities
    $m_0,\ldots,m_{n-1}$ with $\sum_kd_km_k=D$, realized by an explicit
    identification of the index sets, and a unitary $U\in\MN{D}$ such that
    a matrix $X\in\MN{D}$ satisfies $E^*(X)=X$ exactly when
    \begin{align}
        U^\dagger XU
        &=\bigoplus_{k=0}^{n-1}\Id_{m_k}\otimes B_k
        \notag
    \end{align}
    for some matrices $B_k\in\MN{d_k}$; that is, up to reordering the two
    tensor factors of each block,
    \begin{align}
        \Fix(E^*)
        &=U\left(\bigoplus_{k=0}^{n-1}
          \Id_{m_k}\otimes\MN{d_k}\right)U^\dagger.
        \label{eq:spectral_adjoint_fixed_blocks}
    \end{align}
    The positive definite fixed point removes the zero block: this is the
    unital case of the block representation in Equation~(1.39) of
    \cite{Wolf2012Quantum}, invoked by
    \cite[Theorem~6.14]{Wolf2012Quantum}. The general form of
    \cite[Theorem~6.14]{Wolf2012Quantum}, with a zero block and density
    weights $\rho_k$ on the Schr\"odinger-picture fixed points, is not
    asserted here.
\end{theorem}

\begin{proof}
    \leanok
    \uses{thm:adjointFixedPointsStarSubalgebra,
        thm:starSubalgebra_unitary_block_diagonal_iff}
    The fixed points of $E^*$ form a $*$-subalgebra of $\MN{D}$ because
    $E$ is trace-preserving with a positive definite fixed point
    (Theorem~\ref{thm:adjointFixedPointsStarSubalgebra}).
    Theorem~\ref{thm:starSubalgebra_unitary_block_diagonal_iff} supplies
    the unitary $U$, the dimensions and multiplicities, and the equivalence
    between membership and the block form
    (\ref{eq:spectral_adjoint_fixed_blocks}).
\end{proof}

\begin{theorem}[Unitary block form of the fixed points of a unital Kraus map]%
    \label{thm:fixedPoints_block_form}
    \lean{Kraus.fixedPoints_blockDiagonal_iff}
    \leanok
    \uses{def:kraus_fixed_points, def:kraus_adjoint_map, def:unital_kraus}
    Let $E$ be a unital Kraus map on $\MN{D}$ whose adjoint map $E^*$ has a
    positive definite fixed point $\rho>0$. There are $n\in\N$, positive
    dimensions $d_k$ and multiplicities $m_k$ with
    $\sum_kd_km_k=D$, and a unitary $U\in\MN{D}$ such that a matrix
    $X\in\MN{D}$ satisfies $E(X)=X$ exactly when
    \begin{align}
        U^\dagger XU
        &=\bigoplus_{k=0}^{n-1}\Id_{m_k}\otimes B_k
        \notag
    \end{align}
    for some matrices $B_k\in\MN{d_k}$. This is
    Theorem~\ref{thm:adjointFixedPoints_block_form} with the roles of the
    map and its adjoint exchanged.
\end{theorem}

\begin{proof}
    \leanok
    \uses{thm:fixedPointsStarSubalgebra,
        thm:starSubalgebra_unitary_block_diagonal_iff}
    The fixed points of the unital map $E$ form a $*$-subalgebra of
    $\MN{D}$ by Theorem~\ref{thm:fixedPointsStarSubalgebra}, and
    Theorem~\ref{thm:starSubalgebra_unitary_block_diagonal_iff} supplies
    the unitary and the equivalence.
\end{proof}

\begin{definition}[Canonical full-matrix extension of a direct-sum map]%
    \label{def:direct_sum_full_matrix_extension}
    \lean{Matrix.directSumDiagonalEmbedding}
    \lean{Matrix.directSumDiagonalCompression}
    \lean{Matrix.directSumExtension}
    \lean{Matrix.directSumMapExtension}
    \lean{Matrix.directSumTraceAdjointMap}
    \lean{Matrix.directSumTraceAdjointMapBetween}
    \leanok
    Let
    \begin{align}
        \mathcal A
        &=\bigoplus_{k\in I}\MN{n_k}, \notag\\
        \mathcal B
        &=\bigoplus_{l\in J}\MN{m_l},
        \notag
    \end{align}
    and let $\mathcal H_A=\bigoplus_{k\in I}\C^{n_k}$ and
    $\mathcal H_B=\bigoplus_{l\in J}\C^{m_l}$. Write
    $\iota_A,\iota_B$ for the block-diagonal embeddings and
    $\pi_A,\pi_B$ for diagonal-block compression. For a linear map
    $T:\mathcal A\to\mathcal B$, its canonical full-matrix extension is
    \begin{align}
        \widehat T
        &:=\iota_B\circ T\circ\pi_A:
          \End(\mathcal H_A)\longrightarrow\End(\mathcal H_B).
        \label{eq:spectral_direct_sum_extension}
    \end{align}
    For $X\in\mathcal A$ and $Y\in\mathcal B$, the trace adjoint
    $T^*:\mathcal B\to\mathcal A$ is defined by
    \begin{align}
        \sum_{k\in I}\tr(T^*(Y)_kX_k)
        &=\sum_{l\in J}\tr(Y_lT(X)_l).
        \label{eq:spectral_direct_sum_adjoint}
    \end{align}
    This is the block-diagonal realization needed for the fixed-point and
    classification arguments in
    \cite[Appendix~C.4, lines~1980--2003]{Cirac2016MPDO_arXiv}. It is only
    an auxiliary construction: it does not assert the classification
    conclusion of that passage.
\end{definition}

\begin{theorem}[Embedding and compression for the canonical extension]
    \label{thm:direct_sum_extension_embedding_compression}
    \lean{Matrix.directSumDiagonalCompression_embedding}
    \lean{Matrix.directSumExtension_apply}
    \leanok
    \uses{def:direct_sum_full_matrix_extension}
    For every $A\in\mathcal A$, every $X\in\End(\mathcal H_A)$, and every
    linear endomorphism $T:\mathcal A\to\mathcal A$,
    \begin{align}
        \pi_A(\iota_A(A))
        &=A, \notag\\
        \widehat T(X)
        &=\iota_A(T(\pi_A(X))).
        \notag
    \end{align}
\end{theorem}

\begin{proof}\leanok
    The first identity follows by extracting each diagonal block of
    $\iota_A(A)$. The second is the defining evaluation formula for the
    canonical extension.
\end{proof}

\begin{theorem}[Agreement of the two direct-sum extensions]
    \label{thm:direct_sum_map_extension_self}
    \lean{Matrix.directSumMapExtension_self}
    \leanok
    \uses{def:direct_sum_full_matrix_extension}
    If $T:\mathcal A\to\mathcal A$ is an endomorphism, then its full-matrix
    extension as a map between two finite sums agrees with its canonical
    endomorphic extension.
\end{theorem}

\begin{proof}\leanok
    Both maps are $\iota_A\circ T\circ\pi_A$.
\end{proof}

\begin{theorem}[Fixed embedded families under the canonical extension]
    \label{thm:direct_sum_extension_embedding_fixed}
    \lean{Matrix.directSumExtension_embedding_eq_self_iff}
    \leanok
    \uses{def:direct_sum_full_matrix_extension}
    Let $T:\mathcal A\to\mathcal A$ be a linear endomorphism of a finite sum
    of matrix algebras. For every $A\in\mathcal A$,
    \begin{align}
        \widehat T(\iota_A(A))=\iota_A(A)
        &\quad\Longleftrightarrow\quad T(A)=A.
        \notag
    \end{align}
\end{theorem}

\begin{proof}\leanok
    \uses{thm:direct_sum_extension_embedding_compression}
    Diagonal compression is a left inverse of block-diagonal embedding. Hence
    \begin{align}
        \widehat T(\iota_A(A))
        &=\iota_A(T(A)).
        \notag
    \end{align}
    If this equals $\iota_A(A)$, applying diagonal compression gives
    $T(A)=A$. Conversely, substituting $T(A)=A$ gives the fixed-point
    identity.
\end{proof}

Definition~\ref{def:direct_sum_full_matrix_extension} and the trace-adjoint
laws in Lemma~\ref{lem:direct_sum_rectangular_trace_adjoint} provide
supporting facts for the classification argument. By themselves, they do
not prove Schwarz equality or multiplicativity.
Theorem~\ref{thm:direct_sum_inverse_cptp_star_equiv} supplies those two
conclusions; relabeling and dimension matching of the simple summands, and
implementation on each summand by unitary conjugation, remain outside the
present result.

\begin{lemma}[Trace-adjoint laws for two finite sums of matrix algebras]%
    \label{lem:direct_sum_rectangular_trace_adjoint}
    \lean{Matrix.sum_trace_directSumTraceAdjointMapBetween_mul}
    \lean{Matrix.directSumTraceAdjointMapBetween_involutive}
    \lean{Matrix.directSumTraceAdjointMapBetween_id}
    \lean{Matrix.directSumTraceAdjointMapBetween_comp}
    \lean{Matrix.directSumTraceAdjointMapBetween_self}
    \lean{Matrix.IsTracePreservingBetweenDirectSums.directSumTraceAdjointMapBetween_one}
    \lean{Matrix.traceAdjointMap_directSumMapExtension}
    \lean{IsKrausCP.traceAdjointMap}
    \lean{Matrix.IsKrausDirectSumMap.map_posSemidef}
    \lean{Matrix.IsKrausDirectSumMap.directSumTraceAdjointMapBetween}
    \leanok
    \uses{def:direct_sum_full_matrix_extension}
    For maps $T:\mathcal A\to\mathcal B$ and
    $S:\mathcal B\to\mathcal C$, trace adjoints satisfy
    \begin{align}
        \id_{\mathcal A}^*
        &=\id_{\mathcal A}, \notag\\
        (S\circ T)^*
        &=T^*\circ S^*, \notag\\
        (T^*)^*
        &=T, \notag\\
        (\widehat T)^*
        &=\widehat{T^*}.
        \label{eq:spectral_direct_sum_adjoint_laws}
    \end{align}
    If $T$ preserves the total trace, then
    $T^*(\Id_{\mathcal B})=\Id_{\mathcal A}$. If $\widehat T$ admits a
    Kraus representation, then so does $\widehat{T^*}$; in particular, both
    $T$ and $T^*$ send families of positive semidefinite matrices to
    positive semidefinite families. These statements are only the adjoint
    and positivity part of the classification argument; none of its
    multiplicative or blockwise conclusions is asserted here.
\end{lemma}

\begin{proof}
    \leanok
    \uses{def:direct_sum_full_matrix_extension}
    The defining trace identity
    (\ref{eq:spectral_direct_sum_adjoint}) and nondegeneracy of the trace
    pairing give the identity and composition formulas in
    (\ref{eq:spectral_direct_sum_adjoint_laws}). Total-trace preservation
    gives, for every $X\in\mathcal A$,
    \begin{align}
        \sum_{k\in I}\tr(T^*(\Id_{\mathcal B})_kX_k)
        &=\sum_{l\in J}\tr(T(X)_l)
          =\sum_{k\in I}\tr(X_k),
        \notag
    \end{align}
    so $T^*(\Id_{\mathcal B})=\Id_{\mathcal A}$. Finally, if
    $\widehat T(X)=\sum_aK_aXK_a^*$, cyclicity of the trace gives
    $(\widehat T)^*(Y)=\sum_aK_a^*YK_a$. Thus the adjoint again has a
    Kraus representation, and compression of positive block-diagonal
    matrices gives the stated positivity.
\end{proof}

\begin{lemma}[Trace adjoint of a trace-preserving star isomorphism]
    \label{lem:direct_sum_trace_adjoint_star_equiv_inverse}
    \lean{Matrix.directSumTraceAdjointMapBetween_starAlgEquiv_eq_symm}
    \leanok
    Let $\Phi:\mathcal A\simeq\mathcal B$ be a star-algebra isomorphism
    between finite products of full matrix algebras. If $\Phi$ preserves the
    total block trace, then $\Phi^*=\Phi^{-1}$.
\end{lemma}

\begin{proof}
    \leanok
    \uses{lem:direct_sum_rectangular_trace_adjoint}
    For $A\in\mathcal B$ and $B\in\mathcal A$, multiplicativity and trace
    preservation give
    \begin{align}
        \sum_j\tr(A_j\Phi(B)_j)
        &=\sum_j\tr\!\left(\Phi(\Phi^{-1}(A)B)_j\right)
          =\sum_i\tr(\Phi^{-1}(A)_iB_i).
        \notag
    \end{align}
    Nondegeneracy of the trace pairing gives $\Phi^*=\Phi^{-1}$.
\end{proof}

\begin{theorem}[Channel and fixed-point compatibility of the canonical extension]%
    \label{thm:direct_sum_full_matrix_extension_compatibility}
    \lean{Matrix.IsPositiveDirectSumMap.directSumExtension_isPositiveMap}
    \lean{Matrix.IsSchwarzDirectSumMap.directSumExtension_isSchwarzMap}
    \lean{Matrix.IsTracePreservingDirectSumMap.directSumExtension_isTracePreservingMap}
    \lean{Matrix.traceAdjointMap_directSumExtension}
    \lean{Matrix.IsSchwarzDirectSumMap.traceAdjointMap_directSumExtension_isSchwarzMap}
    \lean{Matrix.directSumExtension_apply_eq_self_iff}
    \leanok
    \uses{def:direct_sum_full_matrix_extension}
    In the endomorphic case $\mathcal B=\mathcal A$, write
    $\mathcal H=\mathcal H_A=\mathcal H_B$ and
    $\iota=\iota_A=\iota_B$. If $T$ is positive and satisfies the Schwarz
    inequality on $\mathcal A$, then $\widehat T$ is positive and satisfies
    the Schwarz inequality on $\End(\mathcal H)$. If $T$ preserves the
    total trace $\sum_k\tr(A_k)$, then $\widehat T$ preserves the ordinary
    trace. Moreover, $(\widehat T)^*=\widehat{T^*}$. If $T$ is positive and
    its direct-sum trace adjoint satisfies the Schwarz inequality, then the
    trace adjoint of $\widehat T$ satisfies the Schwarz inequality on
    $\End(\mathcal H)$. Finally,
    \begin{align}
        \Fix(\widehat T)
        &=\iota(\Fix(T)).
        \label{eq:spectral_extension_fixed_points}
    \end{align}
    Equivalently, for every $X\in\End(\mathcal H)$,
    \begin{align}
        \widehat T(X)=X
        &\quad\Longleftrightarrow\quad
        \exists A\in\mathcal A,\quad
        T(A)=A\ \text{ and }\ X=\iota(A).
        \notag
    \end{align}
\end{theorem}

\begin{proof}
    \leanok
    \uses{def:direct_sum_full_matrix_extension}
    For the matrix $R_k$ formed from the column blocks of $A$ outside the
    $k$-th diagonal block, diagonal compression satisfies
    \begin{align}
        \pi(A^*A)_k-\pi(A)_k^*\pi(A)_k
        &=R_k^*R_k\succeq0.
        \notag
    \end{align}
    Block-diagonal embedding is a positive algebra homomorphism and converts
    the sum of the block traces into the full trace. It follows that
    positivity, both stated Schwarz implications, and trace preservation
    pass to $\widehat T$. The embedding and compression satisfy
    \begin{align}
        \tr(\iota(A)X)
        &=\sum_k\tr(A_k\pi(X)_k).
        \notag
    \end{align}
    Substituting (\ref{eq:spectral_direct_sum_extension}) and using
    (\ref{eq:spectral_direct_sum_adjoint}) gives
    $(\widehat T)^*=\widehat{T^*}$.

    Since $\pi\circ\iota=\id$, one fixed-point implication follows from
    \begin{align}
        T(\pi(X))
        &=\pi(\widehat T(X))
          =\pi(X), \notag\\
        X
        &=\widehat T(X)
          =\iota(T(\pi(X)))
          =\iota(\pi(X)),
        \notag
    \end{align}
    whenever $\widehat T(X)=X$. Conversely, if $T(A)=A$, then
    \begin{align}
        \widehat T(\iota(A))
        &=\iota(T(\pi(\iota(A))))
          =\iota(T(A))
          =\iota(A).
        \notag
    \end{align}
    These implications prove
    (\ref{eq:spectral_extension_fixed_points}).
\end{proof}

\section{Unital embeddings of full matrix algebras}

Let $\Phi:\MN{p}\to\MN{q}$ be a unital homomorphism of complex algebras. Its
relative commutant is the algebra
\begin{align}
    \Phi(\MN{p})'
    &=\{T\in\MN{q}\ :\ T\Phi(A)=\Phi(A)T\ \text{for every}\ A\in\MN{p}\},
    \label{eq:spectral_relative_commutant}
\end{align}
which always contains $\C\Id$. This section proves that
$\Phi(\MN{p})'=\C\Id$ forces $\Phi$ to be onto, and records the
contrapositive witness. The argument is the single-summand,
multiplicity-one specialization of the representation argument
of~\cite[Appendix~A]{SchumacherWerner2004}, which decomposes
$\Phi$ into irreducible representations of $\MN{p}$ and uses that all of them
are equivalent to the defining representation on $\C^p$. It is the
finite-dimensional step behind the two support-algebra
inclusions~\cite[Section~7]{GrossNesmeVogtsWerner2012}: a strict inclusion would produce
an element of a relative commutant, or would keep an automorphism from being
onto.

\begin{lemma}[Simplicity of the defining module]%
    \label{lem:matrix_defining_module_simple}
    \lean{Matrix.isSimpleModule_pi}
    \leanok
    Let $K$ be a field and $n$ a nonempty finite index set. Then $K^n$ is a
    simple module over $M_n(K)$.
\end{lemma}

\begin{proof}
    \leanok
    The module is nonzero. Let $x\in K^n$ be nonzero and let $y\in K^n$ be
    arbitrary. Choose $j$ with $x_j\neq0$ and let $A$ be the matrix with
    entries $A_{ik}=\delta_{kj}\,y_i x_j^{-1}$. Then $(Ax)_i=y_ix_j^{-1}x_j=y_i$,
    so the submodule generated by $x$ is everything.
\end{proof}

\begin{lemma}[Uniqueness of the simple module]%
    \label{lem:simple_artinian_simple_modules_isomorphic}
    \lean{IsSimpleRing.nonempty_linearEquiv_of_isSimpleModule}
    \leanok
    Let $R$ be a simple Artinian ring. Any two simple $R$-modules are
    isomorphic.
\end{lemma}

\begin{proof}
    \leanok
    A simple Artinian ring is semisimple, so every simple module over it is
    isomorphic to a simple left ideal, and the regular module of a simple
    Artinian ring is isotypic: its simple left ideals are mutually isomorphic.
    Composing the two isomorphisms with one between the ideals gives the
    claim.
\end{proof}

\begin{theorem}[Equality of matrix sizes]%
    \label{thm:matrix_alg_hom_size_eq_of_scalar_relative_commutant}
    \lean{Matrix.AlgHom.size_eq_of_centralizer_range_eq_bot}
    \leanok
    Let $p,q\geq1$ and let $\Phi:\MN{p}\to\MN{q}$ be a unital homomorphism of
    complex algebras with $\Phi(\MN{p})'=\C\Id$. Then $p=q$.
\end{theorem}

\begin{proof}
    \leanok
    \uses{lem:matrix_defining_module_simple,
        lem:simple_artinian_simple_modules_isomorphic}
    Pulling the matrix action back along $\Phi$ makes $\C^q$ a module over
    $\MN{p}$, with $A\cdot v=\Phi(A)v$; scalars act compatibly because $\Phi$
    is complex-linear. Since $\MN{p}$ is simple Artinian, this module is
    semisimple, so every submodule $W$ has a complement, and the projection
    $\pi$ onto $W$ along that complement is a module map. Read as a $q\times q$
    matrix, $\pi$ satisfies
    \begin{align}
        \pi\Phi(A)v
        &=\pi(A\cdot v)
          =A\cdot\pi v
          =\Phi(A)\pi v,
        \label{eq:spectral_projection_commutes}
    \end{align}
    so $\pi\in\Phi(\MN{p})'=\C\Id$, say $\pi=c\Id$. If $c=0$ then $\pi=0$ and
    $W=0$; if $c\neq0$ then $\pi(c^{-1}v)=v$ for every $v$ and $W=\C^q$. Hence
    $\C^q$ is a simple module over $\MN{p}$. By
    Lemma~\ref{lem:matrix_defining_module_simple} so is the defining module
    $\C^p$, and by
    Lemma~\ref{lem:simple_artinian_simple_modules_isomorphic} the two are
    isomorphic as modules over $\MN{p}$, hence as complex vector spaces.
    Comparing complex dimensions gives $p=q$.
\end{proof}

\begin{theorem}[Scalar relative commutant forces surjectivity]%
    \label{thm:matrix_alg_hom_surjective_of_scalar_relative_commutant}
    \lean{Matrix.AlgHom.surjective_of_centralizer_range_eq_bot}
    \leanok
    Let $p,q\geq1$ and let $\Phi:\MN{p}\to\MN{q}$ be a unital homomorphism of
    complex algebras with $\Phi(\MN{p})'=\C\Id$. Then $\Phi$ is onto.
\end{theorem}

\begin{proof}
    \leanok
    \uses{thm:matrix_alg_hom_size_eq_of_scalar_relative_commutant}
    Injectivity is not assumed: the kernel of $\Phi$ is a two-sided ideal of
    the simple ring $\MN{p}$, and it is not everything because $\MN{q}$ is
    nontrivial, so $\Phi$ is injective. By
    Theorem~\ref{thm:matrix_alg_hom_size_eq_of_scalar_relative_commutant},
    $p=q$, so source and target have the same finite complex dimension $p^2$,
    and an injective linear map between spaces of equal finite dimension is
    onto.
\end{proof}

\begin{corollary}[Star-homomorphism form]%
    \label{cor:matrix_star_alg_hom_surjective_of_scalar_relative_commutant}
    \lean{Matrix.StarAlgHom.surjective_of_centralizer_range_eq_bot}
    \leanok
    Let $p,q\geq1$ and let $\Phi:\MN{p}\to\MN{q}$ be a unital
    $*$-homomorphism with $\Phi(\MN{p})'=\C\Id$. Then $\Phi$ is onto.
\end{corollary}

\begin{proof}
    \leanok
    \uses{thm:matrix_alg_hom_surjective_of_scalar_relative_commutant}
    A unital $*$-homomorphism is in particular a unital homomorphism of
    complex algebras with the same image, so
    Theorem~\ref{thm:matrix_alg_hom_surjective_of_scalar_relative_commutant}
    applies. Preservation of the adjoint plays no part in the argument.
\end{proof}

\begin{theorem}[Non-scalar witness in the relative commutant]%
    \label{thm:matrix_alg_hom_nonscalar_relative_commutant_of_not_surjective}
    \lean{Matrix.AlgHom.exists_commute_notMem_bot_of_not_surjective}
    \leanok
    Let $p,q\geq1$ and let $\Phi:\MN{p}\to\MN{q}$ be a unital homomorphism of
    complex algebras that is not onto. Then there is a $T\in\MN{q}$ commuting
    with every $\Phi(A)$ and not a scalar multiple of $\Id$.
\end{theorem}

\begin{proof}
    \leanok
    \uses{thm:matrix_alg_hom_surjective_of_scalar_relative_commutant}
    If no such $T$ existed, then every element of the relative commutant
    (\ref{eq:spectral_relative_commutant}) would be a scalar multiple of
    $\Id$, so $\Phi(\MN{p})'=\C\Id$ and
    Theorem~\ref{thm:matrix_alg_hom_surjective_of_scalar_relative_commutant}
    would make $\Phi$ onto.
\end{proof}
